% T20 — «Бланк ЕГЭ»: гриф, водяной знак «ТРЕНАЖЁР» по диагонали, поле ФИО, вариант.
% Профиль: математика 11 класс, 7 задач, 1 кол.
\documentclass[a4paper,11pt]{article}

\usepackage{../../../../Lessons/_templates/preamble_worksheet}
\usepackage{../_shared/worksheet_check}

\usepackage{eso-pic}

% --- Палитра: серый + чёрный, без эмоций ---
\definecolor{wcprimary}{HTML}{2D3748}
\definecolor{wcprimaryLight}{HTML}{CBD5E0}
\definecolor{wcprimaryBG}{HTML}{F7FAFC}
\definecolor{examgray}{HTML}{4A5568}
\definecolor{examwm}{HTML}{E2E8F0}

% Засечный (по умолчанию из preamble — Times)
\renewcommand{\familydefault}{\rmdefault}

\geometry{a4paper, top=20mm, bottom=20mm, left=18mm, right=18mm}

% --- Водяной знак «ТРЕНАЖЁР» по диагонали через eso-pic ---
\AddToShipoutPictureBG{%
  \begin{tikzpicture}[remember picture, overlay]
    \node[rotate=45, anchor=center, font=\sffamily\bfseries\fontsize{72pt}{72pt}\selectfont,
          text=examwm, opacity=0.45]
      at (current page.center) {ТРЕНАЖЁР};
  \end{tikzpicture}}

% --- Заголовок: официальный гриф, бланк ---
\renewcommand{\WorksheetTitle}[1]{%
  \par\noindent\begin{tikzpicture}
    % Верхняя серая полоска
    \fill[wcprimary] (0,0) rectangle (\linewidth, -2mm);
    \node[anchor=north west, font=\sffamily\bfseries\fontsize{9pt}{11pt}\selectfont,
          text=wcprimary] at (0, -4mm)
      {ТРЕНИРОВОЧНАЯ РАБОТА В ФОРМАТЕ ЕГЭ \ \ \textbar\ \ ПРОФИЛЬНЫЙ УРОВЕНЬ};
    \node[anchor=north east, font=\sffamily\fontsize{9pt}{11pt}\selectfont,
          text=examgray] at (\linewidth, -4mm)
      {Вариант \No 1};
  \end{tikzpicture}\par
  \vspace{5mm}
  \noindent{\fontsize{18pt}{22pt}\selectfont\bfseries\color{wcprimary}#1}\par
  \vspace{1mm}{\color{wcprimary}\rule{\linewidth}{0.6pt}}\par
  \vspace{4mm}}

% --- TaskBox: ¹Задание N — узкая шапка, тонкая линия снизу ---
\renewcommand{\TaskBox}[2]{%
  \par\nopagebreak\noindent
  \begin{tikzpicture}
    \node[fill=wcprimary, text=white, font=\sffamily\bfseries\footnotesize,
          inner xsep=3mm, inner ysep=1mm, anchor=west] (T) at (0,0)
      {ЗАДАНИЕ \No #1};
  \end{tikzpicture}\par\vspace{1.5mm}
  \noindent #2\par
  \vspace{1mm}{\color{wcprimary!40}\rule{\linewidth}{0.3pt}}\par\vspace{3mm}}

\SetAnswerFieldSize{34mm}{8mm}

% Кастомный header с ФИО+класс+школа+дата
\newcommand{\ExamFormHeader}{%
  \noindent\begin{tikzpicture}
    \node[draw=wcprimary, line width=0.8pt, inner sep=2.5mm,
          text width=\linewidth, font=\sffamily\small] {%
      \textbf{Фамилия:}\ \rule{0.32\linewidth}{0.4pt}\hfill
      \textbf{Имя:}\ \rule{0.22\linewidth}{0.4pt}\hfill
      \textbf{Отчество:}\ \rule{0.22\linewidth}{0.4pt}\\[2mm]
      \textbf{Школа:}\ \rule{0.40\linewidth}{0.4pt}\hfill
      \textbf{Класс:}\ \rule{0.08\linewidth}{0.4pt}\hfill
      \textbf{Дата:}\ \rule{0.15\linewidth}{0.4pt}};
  \end{tikzpicture}\par\vspace{4mm}}

\begin{document}

\WorksheetTitle{Контрольно-измерительные материалы}
\par\noindent{\sffamily\small\color{examgray}Демонстрационный вариант \quad·\quad класс 11 \quad·\quad T20 \quad·\quad ID: T20-exam-2026-05-27}\par\vspace{3mm}
\ExamFormHeader

\TaskBox{1}{%
  Найдите значение выражения $\dfrac{14^{2{,}7}}{14^{0{,}7}}$.\par
  \vspace{1mm}\hfill\textbf{Ответ:}\ \AnswerField{1}%
}

\TaskBox{2}{%
  Поезд, двигаясь равномерно со скоростью $60$~км/ч, проезжает мимо стоящего у обочины пешехода за $20$~секунд. Найдите длину поезда в метрах.\par
  \vspace{1mm}\hfill\textbf{Ответ:}\ \AnswerField{2}%
}

\TaskBox{3}{%
  В равнобедренном треугольнике $ABC$ боковая сторона $AB = 10$, основание $AC = 12$. Найдите радиус вписанной окружности.\par
  \vspace{1mm}\hfill\textbf{Ответ:}\ \AnswerField{3}%
}

\TaskBox{4}{%
  Решите уравнение $\log_3(x - 2) + \log_3(x + 6) = 2$.\par
  \vspace{1mm}\hfill\textbf{Ответ:}\ \AnswerField{4}%
}

\TaskBox{5}{%
  Найдите наименьшее значение функции $y = e^{2x} - 4e^x + 7$.\par
  \vspace{1mm}\hfill\textbf{Ответ:}\ \AnswerField{5}%
}

\TaskBox{6}{%
  Найдите площадь треугольника со сторонами $5$, $12$ и $13$.\par
  \vspace{1mm}\hfill\textbf{Ответ:}\ \AnswerField{6}%
}

\TaskBox{7}{%
  Вероятность того, что батарейка бракованная, равна $0{,}05$. Найдите вероятность, что среди двух взятых батареек обе исправны. Округлите до сотых.\par
  \vspace{1mm}\hfill\textbf{Ответ:}\ \AnswerField{7}%
}

\WorksheetQR{T20-exam/L1/v1}

\end{document}
